\clearpage
\VerdierPhysicalPageAnchor{0067}%
\VerdierSemanticAnchor{ch1-2-1-5}
\noindent\textbf{2.1.5.}\quad
射 $f^\sigma\colon Y^\sigma\to Y'^\sigma$ を射 $(f^\sigma)$ の
\emph{次数 $\sigma$ の成分}と呼ぶ。$f\colon Y\to Y'$ を $n$-重複複体の
射とするとき、$f$ の次数 $\sigma$ の成分を $f^\sigma$ と記す。
\VerdierSemanticAnchor{ch1-2-1-5-composition}
$n$-重複複体の射は明らかな仕方で、成分ごとに合成される。
\VerdierSemanticAnchor{ch1-2-1-5-category}
$C$ の対象からなる $n$-重複複体は、$\operatorname{comp}^{n}(C)$ と記す
圏をなす。$n=1$ のときは簡略記法
$\operatorname{comp}^{1}(C)=\operatorname{comp}(C)$ を用いる。
$\operatorname{comp}(C)$ の対象を $C$ における\emph{単純複体}、また混同の
おそれがなければ単に $C$ における\emph{複体}と呼ぶ。圏
$\operatorname{comp}^{n}(C)$ は加法圏である。

\VerdierSemanticAnchor{ch1-2-1-6}
\noindent\textbf{2.1.6.}\quad
$Y$ を $C$ における $n$-重複複体とする。$Y$ の次数 $\sigma$ の成分を
$Y_{-\sigma}$、方向 $i$ における次数 $\sigma$ の微分を
$d^Y_{-\sigma,i}$ と記すこともある。
\VerdierSemanticAnchor{ch1-2-1-6-positive-finite}
この記法は、とくに $Y$ が単純複体で、正の次数の成分が有限個を除いて
消える場合に用いる。同様に、$f\colon Y\to Y'$ を $n$-重複複体の射と
すると、射 $f$ の次数 $\sigma$ の成分を $f_{-\sigma}$ と記すことがある。

\VerdierSemanticAnchor{ch1-2-1-7}
\noindent\textbf{2.1.7.}\quad
$\varphi\colon[n]\to[m]$ を写像とする。
\VerdierSemanticAnchor{ch1-2-1-7-summable}
各 $\tau\in\mathbb Z[m]$ に対して、$\varphi(\sigma)=\tau$ かつ
$Y^\sigma\neq0$ となる $\sigma\in\mathbb Z[n]$ が有限個しかないとき、
$C$ における $n$-重複複体 $Y$ を\emph{$\varphi$-総和可能}という。
$\varphi$ が唯一の写像 $[n]\to[1]$ であるときは、単に\emph{総和可能}という。
\VerdierSemanticAnchor{ch1-2-1-7-composition}
$[n]\xrightarrow{\varphi}[m]\xrightarrow{\psi}[p]$ を二つの写像とする。
$\psi\varphi$-総和可能な任意の $n$-重複複体 $Y$ は $\varphi$-総和可能でも
ある。とくに、総和可能な任意の $n$-重複複体は、任意の写像
$[n]\xrightarrow{\varphi}[m]$ に対して $\varphi$-総和可能である。
\VerdierSemanticAnchor{ch1-2-1-7-injective}
$\varphi$ が単射ならば、すべての複体は $\varphi$-総和可能である。
\VerdierSemanticAnchor{ch1-2-1-7-category}
\emph{$\varphi$-総和可能な $n$-重複複体の圏}とは、
$\varphi$-総和可能な複体によって定まる $\operatorname{comp}^{n}(C)$ の
充満部分圏である。

\VerdierSemanticAnchor{ch1-2-2}
\subsection{付随する単純複体}

\VerdierSemanticAnchor{ch1-2-2-1}
\noindent\textbf{2.2.1.}\quad
各 $\sigma\in\mathbb Z[n]$ を
$\sigma=\sum_{j\in[n]}l_je_j$、$l_j\in\mathbb Z$ と書き、各
$i\in[n]$ に対して

\VerdierSemanticAnchor{ch1-2-2-1-eq1}
\begin{equation}
  \varepsilon(\sigma,i)=(-1)^{\left(\sum_{j<i}l_j\right)}
  \tag{2.2.1.1}
\end{equation}
とおく。

\VerdierSemanticAnchor{ch1-2-2-1-sign-map}
公式 (2.2.1.1) は $\mathbb Z[n]\times[n]$ から符号群 $\{-1,1\}$ への写像
\begin{equation*}
  (\sigma,i)\longmapsto\varepsilon(\sigma,i)
\end{equation*}
を定める。各 $i\in[n]$ について
$\sigma\longmapsto\varepsilon(\sigma,i)$ が群準同型であり、
